\DocumentMetadata{
  pdfversion=2.0,
  pdfstandard=ua-2,
  lang=en-US,
  tagging=on,
  tagging-setup={math/setup=mathml-SE,tabsorder=structure}
}
\documentclass[11pt,letterpaper]{article}
\usepackage[margin=0.68in,includefoot]{geometry}
\usepackage{newcomputermodern}
\usepackage{amsmath,amscd,mathtools}
\usepackage{tikz,adjustbox}
\providecommand{\square}{\mdlgwhtsquare}
\usepackage[hidelinks]{hyperref}
\hypersetup{
  pdftitle={Algebra PhD Exam, August 2022},
  pdfauthor={University of Florida Department of Mathematics},
  pdfsubject={Graduate mathematics examination},
  pdfdisplaydoctitle=true,
  bookmarksopen=true
}
\setcounter{secnumdepth}{0}
\setlength{\parindent}{0pt}
\setlength{\parskip}{0.28em}
\setlength{\emergencystretch}{3em}
\setlength{\leftmargini}{2.15em}
\setlength{\leftmarginii}{2.65em}
\setlength{\leftmarginiii}{3em}
\pagestyle{empty}
\raggedbottom
\allowdisplaybreaks
\NewTaggingSocketPlug{block/list/label}{examproblem}{%
  \tagstructbegin{tag=Lbl}\tagstructend%
  \tagstructbegin{tag=\LBody}%
  \tagstructbegin{tag=\ProblemHeadingTag,title-o={\ProblemHeadingText}}%
  \tagmcbegin{tag=\ProblemHeadingTag,actualtext={\ProblemHeadingText}}%
  #2%
  \tagmcend%
  \tagstructend%
}
\newenvironment{examproblems}[2]{%
  \def\ProblemHeadingTag{#2}%
  \AssignTaggingSocketPlug{block/list/label}{examproblem}%
  \begin{enumerate}%
  \setcounter{enumi}{#1}%
  \renewcommand{\labelenumi}{\textbf{\theenumi.}}%
  \setlength{\topsep}{0.35em}%
  \setlength{\partopsep}{0pt}%
  \setlength{\itemsep}{0.48em}%
  \setlength{\parsep}{0.18em}%
}{\end{enumerate}}
\newenvironment{examsubparts}{%
  \AssignTaggingSocketPlug{block/list/label}{default}%
  \begin{enumerate}%
  \setlength{\topsep}{0.18em}%
  \setlength{\partopsep}{0pt}%
  \setlength{\itemsep}{0.16em}%
  \setlength{\parsep}{0.1em}%
}{\end{enumerate}}
\newenvironment{examinstructions}{%
  \par\begingroup\itshape
}{%
  \par\endgroup\vspace{0.18em}
}
\makeatletter
\renewcommand\subsection{\@startsection{subsection}{2}{0pt}%
  {-1.25ex plus -.25ex minus -.1ex}%
  {0.55ex plus .1ex}%
  {\normalfont\large\bfseries}}
\renewcommand{\@maketitle}{%
  \newpage\null\vskip -1em%
  {\footnotesize\bfseries
    \href{https://www.ufl.edu/}{University of Florida}, \href{https://math.ufl.edu/}{Department of Mathematics}\par}%
  \vskip -0.5em%
  \tagstructbegin{tag=H1,title={Algebra PhD Exam, August 2022}}%
  \tagpdfparaOff%
  \tagmcbegin{tag=H1}%
    {\LARGE\bfseries\href{https://gma.math.ufl.edu/exams/algebra-phd/}{Algebra PhD Exam}, August 2022\par}%
  \tagmcend%
  \tagpdfparaOn%
  \tagstructend%
  \vskip 0.25em%
  \tagmcbegin{artifact}%
    \hrule height 1pt%
  \tagmcend%
  \vskip 1em%
}
\makeatother
\title{Algebra PhD Exam, August 2022}
\author{}
\date{}
\begin{document}
\maketitle
\thispagestyle{empty}
\begin{examinstructions}
Answer seven problems. Write your answers clearly in complete English sentences. You may quote results (within reason) as long as you state them clearly.
\end{examinstructions}
\begin{examproblems}{0}{H2}
\def\ProblemHeadingText{Problem 1.}
\item
This problem has two parts.
\begin{examsubparts}
\item[\textnormal{(a)}]
Give the definition of a free object in a concrete category \(\mathcal C\).
\item[\textnormal{(b)}]
Let \(\mathcal G\) be the category of finite groups. Prove that an object \(H \in \mathcal G\) is free if and only if~\(H\) is a trivial group.
\end{examsubparts}
\def\ProblemHeadingText{Problem 2.}
\item
This problem has three parts.
\begin{examsubparts}
\item[\textnormal{(a)}]
Define what it means for a group~\(G\) to be solvable.
\item[\textnormal{(b)}]
Let~\(G\) be a solvable group and let~\(N\) be a normal subgroup of~\(G\). Prove that \(G/N\) is solvable.
\item[\textnormal{(c)}]
Let~\(F\) be a group which is free on \(n>1\) generators. Prove that~\(F\) is not solvable.
\end{examsubparts}
\def\ProblemHeadingText{Problem 3.}
\item
Let \(m,n\) be positive integers. Find~\(r\) such that there is an isomorphism of \(\mathbb Z\)-modules
\[
(\mathbb Z/m\mathbb Z)\otimes_{\mathbb Z}(\mathbb Z/n\mathbb Z)\cong \mathbb Z/r\mathbb Z.
\]
 Justify your answer.
\def\ProblemHeadingText{Problem 4.}
\item
Let~\(R\) be a ring. Prove that the following are equivalent:
\begin{examsubparts}
\item[\textnormal{(a)}]
Every~\(R\)-module is projective.
\item[\textnormal{(b)}]
Every~\(R\)-module is injective.
\item[\textnormal{(c)}]
Every short exact sequence of~\(R\)-modules splits.
\end{examsubparts}
\def\ProblemHeadingText{Problem 5.}
\item
Let~\(G\) be a group, let~\(F\) be a field, and let \(\chi_1,\ldots,\chi_n\) be distinct homomorphisms from~\(G\) into \(F^\times\). Prove that \(\chi_1,\ldots,\chi_n\) are linearly independent over~\(F\).
\def\ProblemHeadingText{Problem 6.}
\item
This problem has two parts.
\begin{examsubparts}
\item[\textnormal{(a)}]
Let \(n\geq 1\). Define the~\(n\)th cyclotomic polynomial \(\Phi_n(X)\).
\item[\textnormal{(b)}]
Prove that \(\Phi_n(X)\) is irreducible over \(\mathbb Q\). You may assume that \(\Phi_n(X)\in\mathbb Z[X]\).
\end{examsubparts}
\def\ProblemHeadingText{Problem 7.}
\item
Let~\(R\) be a commutative Noetherian ring with~\(1\).
\begin{examsubparts}
\item[\textnormal{(a)}]
Define what it means for a proper ideal \(J\subset R\) to be irreducible.
\item[\textnormal{(b)}]
Prove that every irreducible proper ideal in~\(R\) is primary.
\item[\textnormal{(c)}]
Let \(I\subset R\) be an ideal. Prove that~\(I\) has a primary decomposition.
\end{examsubparts}
\def\ProblemHeadingText{Problem 8.}
\item
Let~\(S\) be a commutative ring with~\(1\), let~\(R\) be a subring of~\(S\) which contains~\(1\), and let \(x\in S\). Suppose there is a subring \(T\subset S\) such that \(T\supset R[x]\) and~\(T\) is finitely generated as an~\(R\)-module. Prove that~\(x\) is a root of some monic polynomial with coefficients in~\(R\).
\def\ProblemHeadingText{Problem 9.}
\item
Let~\(R\) be a Noetherian integral domain with a unique nonzero prime ideal~\(P\). Assume that~\(R\) is integrally closed in its fraction field~\(F\). Prove that~\(P\) is a principal ideal.
\def\ProblemHeadingText{Problem 10.}
\item
Determine which of the following are Dedekind domains. Give a brief explanation in each case.
\begin{examsubparts}
\item[\textnormal{(a)}]
\(\mathbb Z[\sqrt{3}]\)
\item[\textnormal{(b)}]
\(\mathbb Z[\sqrt{5}]\)
\item[\textnormal{(c)}]
\(\mathbb Q[X]\)
\item[\textnormal{(d)}]
\(\mathbb Q[X,Y]\)
\end{examsubparts}
\def\ProblemHeadingText{Problem 11.}
\item
Let~\(D\) be a finite division ring. Prove that~\(D\) is a field.
\end{examproblems}
\end{document}
